\documentclass{beamer}
\usepackage{graphicx}
\usepackage{listings} % Syntax highlighing
\usepackage{fancyvrb} % Inline verbatim
\usepackage{hyperref} % Hyperlinks
\hypersetup{pdfpagemode=FullScreen}
\usepackage{tikz}

\usetheme{Boadilla}
\title{Odds \& Ends}
\author{UMBC Malware Data Science}
\date{Week 6: 3 March 2020}

\begin{document}

\begin{frame}{Data Science \& Cyber}
    Data Science: a discipline which uses data to answer tough questions; more data generally yields better answers.
    \\ ~~ \\
    Cybersecurity: a discipline with a lot of tough questions needing answering.
\end{frame}

\begin{frame}{Conferences}
    Noteworthy conferences:
    
    \begin{itemize}
        \item \href{https://www.shmoocon.org/}{Shmoocon \includegraphics[scale=0.02]{Images/ShmooEyes.jpg}}: DC, January
        \item \href{https://bsidescharm.org/}{BSides Charm}: Baltimore, Spring
        \item \href{https://www.nvidia.com/en-us/gtc/}{GTC \includegraphics[scale=0.15]{Images/gtc.png}}: San Jose, Spring; DC, Fall
        \item \href{https://twitter.com/hashtag/hackersummercamp?src=hashtag_click}{\#HackerSummerCamp}:
        \begin{itemize}
            \item \href{https://www.blackhat.com/}{Blackhat \includegraphics[scale=0.025]{Images/blackhat.png}}: Las Vegas, Early August
            \item \href{https://www.bsideslv.org/}{BSides Las Vegas}: same time
            \item \href{https://defcon.org/}{Defcon\includegraphics[scale=0.06]{Images/Defcon.png}}: Las Vegas, after Black Hat
        \end{itemize}
        \item \href{https://hope.net/}{HOPE}: NYC, Summer
        \item \href{https://toorcon.net/}{toorCon}: Summer
        \item \href{http://www.bsidesdc.org/}{BSides DC}: DC, Fall
        \item \href{https://www.camlis.org/}{CAMLIS}: DC, Fall
        \item \href{http://www.jailbreaksecuritysummit.com/}{Jailbreak Sec}: Laurel, Fall
        \item \href{https://www.malwareconference.org/}{IEEE MalCon}: Nantucket, Fall
    \end{itemize}
        
\end{frame}

\begin{frame}{Cyber Data Science}
    There is an overlap in cybersecurity and data science:
    \begin{itemize}
        \item Elastic bought Endgame in 2019
        \begin{itemize}
            \item Elastic has products for data analytics, big data, and log processing
            \item With Endgame, they now have end point security \& malware analysis expertise
            \item Endgame's open source project \href{https://github.com/endgameinc/eql}{Event Query Language} allows SQL-like queries over event logs
        \end{itemize}
        \item Splunk also has end-to-end log ingest, storage, analysis, retrieval, ML algorithms
        \item Jobs ranging from \$80k to \$200k according to \href{https://www.glassdoor.com/Job/cyber-security-data-scientist-jobs-SRCH_KO0,29.htm}{Glassdoor.com}.
    \end{itemize}
\end{frame}

\begin{frame}{Future Threats}
    As cybersecurity challenges evolve, machine learning might be able to help:
    \begin{itemize}
        \item Deep Fakes
        \begin{itemize}
            \item Phone number \& voice match the CEO, is he really asking you to move money?
            \item Fake video or Instagram post from a celebrity trashes a brand, resulting in stock price plummet
        \end{itemize}
        \item Ransomware $\rightarrow$ Extortionware
        \begin{itemize}
            \item Pay to get your files back,
            \item Pay again or else everyone gets them too!
        \end{itemize}
        \item 5G \& IoT
        \begin{itemize}
            \item IoT is notoriously insecure.
            \item With 5G, more insecure devices will be online.
            \item 5G inherits some 4G vulnerabilities which went unpatched
        \end{itemize}
    \end{itemize}
    Read more: \url{https://www.zdnet.com/article/2020-is-when-cybersecurity-gets-even-weirder-so-get-ready/}
\end{frame}

\begin{frame}{Hardware AI Accelerators}
    \begin{itemize}
        \item GPUs: \href{https://www.nvidia.com/en-us/shop/geforce/}{Nvidia} \href{https://developer.nvidia.com/cuda-zone}{CUDA}
        \item Intel: \href{https://www.movidius.com/}{Movidius}, \href{https://www.intel.ai/nervana-nnp/}{Nervana}
        \item Google: \href{https://cloud.google.com/tpu/}{Tensor Processing Unit (TPU)}, \href{https://coral.ai/products/dev-board/}{Google Coral}
        \item Apple: \href{https://en.wikipedia.org/wiki/Apple_A11}{Neural Engine}
        \item Tesla: \href{https://www.anandtech.com/show/14766/hot-chips-31-live-blogs-tesla-solution-for-full-self-driving}{AP3}
        \item IBM: \href{https://www.ibm.com/blogs/research/tag/truenorth/}{TrueNorth}
        \item Cognimem: CM1K
    \end{itemize}
    
    \\ ~~ \\
    A demonstration of why GPUs perform well for some operations: \url{https://youtu.be/-P28LKWTzrI}.
\end{frame}

\begin{frame}{GPU Programming}
    Two standards:
    \begin{itemize}
        \item CUDA: Made by Nvidia, requires Nvidia GPU, drivers, and toolkit. Nvidia also has specialized software packages for basic applications, such as machine learning, signal processing, matrix multiplication, etc.
        \item OpenCL: Loose standard, works on Nvidia GPUs, AMD GPUs, newer Intel CPUs, some FPGAs, etc.
    \end{itemize}
    \\ ~~ \\
    ML frameworks for GPUs:
    \begin{itemize}
        \item PyTorch
        \item Keras
        \item TensorFlow
        \item Others
    \end{itemize}
\end{frame}

\begin{frame}[fragile]{CUDA Example}
    \lstset{language=C,
    basicstyle=\tiny,
    keywordstyle=\color{blue}\ttfamily,
    stringstyle=\color{red}\ttfamily,
    commentstyle=\color{green}\ttfamily,
    morecomment=[l][\color{magenta}]{\#}
}
\begin{lstlisting}
// Using CUDA device to calculate pi
#include <stdio.h>
#include <cuda.h>

#define NBIN  10000000  // Number of bins
#define NUM_BLOCK   13  // Number of thread blocks
#define NUM_THREAD 192  // Number of threads per block

// Kernel that executes on the CUDA device
__global__ void cal_pi(float *sum, int nbin, float step, int nthreads, int nblocks) {
	int i;
	float x;
	int idx = blockIdx.x*blockDim.x+threadIdx.x;  // Sequential thread index 
	                                              // across the blocks
	for (i=idx; i< nbin; i+=nthreads*nblocks) {
		x = (i+0.5)*step;
		sum[idx] += 4.0/(1.0+x*x);
	}
}
\end{lstlisting}
\end{frame}

\begin{frame}[fragile]{CUDA Example Continued}
    \lstset{language=C,
    basicstyle=\tiny,
    keywordstyle=\color{blue}\ttfamily,
    stringstyle=\color{red}\ttfamily,
    commentstyle=\color{green}\ttfamily,
    morecomment=[l][\color{magenta}]{\#}
}
\begin{lstlisting}
// Using CUDA device to calculate pi
#include <stdio.h>
#include <cuda.h>

#define NBIN  10000000  // Number of bins
#define NUM_BLOCK   13  // Number of thread blocks
#define NUM_THREAD 192  // Number of threads per block

int main(void) {
	dim3 dimGrid(NUM_BLOCK,1,1);  // Grid dimensions
	dim3 dimBlock(NUM_THREAD,1,1);  // Block dimensions
	float *sumHost, *sumDev;  // Pointer to host & device arrays

	float step = 1.0/NBIN, pi = 0;  // Step size
	size_t size = NUM_BLOCK*NUM_THREAD*sizeof(float);  //Array memory size
	sumHost = (float *)malloc(size);  //  Allocate array on host
	cudaMalloc((void **) &sumDev, size);  // Allocate array on device
	// Initialize array in device to 0
	cudaMemset(sumDev, 0, size);
	// Do calculation on device, call CUDA kernel
	cal_pi <<<dimGrid, dimBlock>>> (sumDev, NBIN, step, NUM_THREAD, NUM_BLOCK);
	// Retrieve result from device and store it in host array
	cudaMemcpy(sumHost, sumDev, size, cudaMemcpyDeviceToHost);
	for(int tid=0; tid<NUM_THREAD*NUM_BLOCK; tid++)
		pi += sumHost[tid];
	pi *= step;
	// Print results
	printf("PI = %f\n",pi);
	// Cleanup
	free(sumHost); 
	cudaFree(sumDev);
	return 0;
}
\end{lstlisting}
\end{frame}

\begin{frame}{Unpackers}
    As we know, packed files are hard to analyze statically. Some packers are known, some are readily accessible, some are custom. Wouldn't it be nice to have a tool to automatically unpack malware?
    \pause
    \\ ~~ \\
    This amazing tool doesn't exist.
    \pause
    \\ ~~ \\
    But maybe we can build one...
\end{frame}

\begin{frame}{Binee}
    Check out Binee: \url{https://github.com/carbonblack/binee}.
    \begin{itemize}
        \item It kinda executes software and steps though system calls, but doesn't run the program. It uses \texttt{libunicorn} to emulate an Intel CPU, and keeps track of memory the program would be using if it were running.
        \item Binee also works on packed files.
        \item Therefore Binee is almost a universal unpacker!
    \end{itemize}
    \\ ~~ \\
    We'd just need to study the Binee source code (it's freely available), and figure out where to stop the process, figure out where in the Binee variables the memory information is saved, dump that memory to disk.
    \\ ~~ \\
    This is likely far more difficult to do, but would be very useful.
\end{frame}

\begin{frame}{AVClass}
    \begin{itemize}
        \item \href{https://www.virustotal.com/gui/file/00a0a07947c02a2ebf9f703224ad76755da761d25781620aa26cf5cf14e92b19/detection}{Most malware names from anti-virus products aren't helpful}. 
        \begin{itemize}
            \item Win32:Malware-gen
            \item Generic7.CGLT
            \item Gen:Variant.Symmi.69954
        \end{itemize}
        \item AVClass uses common output from several AVs to determine a malware family name
        \item In the VMs, the Ember 2018 dataset has some malware organized by AVClass malware names
        \item Allows you to better identify which malware might be under-presented in training, though a good model should identify malware families it hasn't seen.
    \end{itemize}
    \footnotesize
    Code: \url{https://github.com/malicialab/avclass} \\
    \footnotesize
    Paper: \url{https://software.imdea.org/~juanca/papers/avclass_raid16.pdf}
\end{frame}

%\begin{frame}{\LaTeX}
%    \LaTeX is a typesetting language, which creates PDF documents. It has a variety of advantages over MS Word, Libreoffice.
%    \begin{itemize}
%        \item It can create a variety of document types: reports, books, articles, posters, and even presentations (like this one).
%        \item Raw data can be used to create charts.
%        \item It's the language often used when submitting articles to journals, conferences.
%        \item The free online \LaTeX service \href{https://www.overleaf.com}{Overleaf.com} makes using it easy, and has tutorials to help.
%    \end{itemize}
%\end{frame}

\begin{frame}{Data Pre-Processing}
    By changing how the data is presented to the model, we can sometimes help the model learn the concept quicker.
    \begin{itemize}
        \item Categories \& strings become numbers
        \begin{itemize}
            \item 32-bit vs. 64-bit could become \texttt{0} or \texttt{1}
            \item Intel, ARM, SPARC, MIPS could become: \texttt{0}, \texttt{1}, \texttt{3}, or \texttt{4}
        \end{itemize}
        \item Scale a range of possible values to [0.0, 1.0]
        \begin{itemize}
            \item Entropy: 4.03 $\rightarrow$ 0.5037
        \end{itemize}
        \item Both for features which don't have a well-defined range, like file size.
        \begin{itemize}
            \item $<1024$ $\rightarrow$ \texttt{0}
            \item 1 kb to 1 MB $\rightarrow$ \texttt{1}
            \item 1 MB to 10 MB $\rightarrow$ \texttt{2}
            \item $>10$ MB $\rightarrow$ \texttt{3}
        \end{itemize}
    \end{itemize}
\end{frame}

\begin{frame}{Lab}
    Explore data pre-processing functionality of \texttt{sklearn}.
    \\ ~~ \\
    Presentations for Project 1 next week!
\end{frame}

\end{document}